\begin{problem}
    Una persona necesita comprar \SI{3}{kg} de azúcar, \SI{6}{L} de lavaozza, \SI{4}{L} de leche y \SI{1}{kg} de arroz.

    En la siguiente tabla se registraron los precios de los productos por comprar.

    \[
    \begin{array}{c|c|c}
    \text{Producto} & \text{Precio unitario} & \text{Promoción por dos unidades} \\
    \hline
    1 \text{ kilogramo de azúcar} & \SI{1400}[\$]{} & \SI{2600}[\$]{} \\
    \frac{1}{2} \text{ litro de lavaozza} & \SI{1600}[\$]{} & \SI{3000}[\$]{} \\
    1 \text{ litro de leche} & \SI{1200}[\$]{} & \SI{2200}[\$]{} \\
    1 \text{ kilogramo de arroz} & \SI{1700}[\$]{} & \SI{3300}[\$]{}
    \end{array}
    \]

    Para calcular el menor precio por pagar, la persona efectúa el siguiente procedimiento, en el cual comete un error.

    \begin{itemize}
        \item Paso 1: calcula que en azúcar debe gastar $\SI{2600}[\$]{} + \SI{1400}[\$]{} = \SI{4000}[\$]{}$.
        \item Paso 2: calcula que en lavaozza debe gastar $3 \cdot \SI{3000}[\$]{} = \SI{9000}[\$]{}$.
        \item Paso 3: calcula que en leche debe gastar $2 \cdot \SI{2200}[\$]{} = \SI{4400}[\$]{}$.
        \item Paso 4: suma a los valores anteriores el precio de un kilogramo de arroz y obtiene que el precio final por pagar es \SI{19100}[\$]{}.
    \end{itemize}

    ¿En cuál de los pasos la persona cometió el error?

    \begin{enumerate}[label=\Alph*.]
        \item En el Paso 1
        \item En el Paso 2
        \item En el Paso 3
        \item En el Paso 4
    \end{enumerate}
\end{problem}
